\subsection*{Introit: a holy place made of concord\enspace\properrefs{Int.}}

The antiphon is not a continuous biblical sentence. It joins the close of
Psalm~67:6, the opening of v.~7, and a clause from v.~36, then supplies v.~2
as the psalm verse. The omissions matter. The complete poem praises the God
who fathers orphans, judges widows, brings prisoners out in strength, leads
Israel from Sinai, takes up residence in Zion, and finally gives strength to
his people. The selected clauses compress that history into dwelling,
concord, strength, and the scattering of enemies.

One word also preserves a different textual history. The Missal sings
\latin{unánimes in domo}; the Clementine has \latin{unius moris in domo},
and the Douay accordingly says ``men of one manner.'' Cassiodorus's lemma has
the chant's \latin{unanimes}. Hilary reads a closely related Old Latin form,
and Augustine glosses his \latin{unius modi} as persons of one mind and one
understanding. The agreement lies in the exegesis, not in identical biblical
wording. This is why the appointed English prints complete Douay verses and
does not invent an English centonization.

Hilary, \emph{Tractatus super Psalmos} 67.8, joins the house to Acts~4:32,
the believers of one heart and soul. Augustine, \emph{Enarratio in Psalmum}
67.7--8, identifies the holy place with persons whom grace builds into
concord; God is not confined to space, and therefore the search for his holy
place turns toward a people made fit for his indwelling. Cassiodorus likewise
names the Church and cites Acts~4:32 and Psalm~132:1. Theodoret reads the
Greek \emph{monotropous} more ascetically as persons directed toward one aim,
virtue. These are real developments: ecclesial concord in Hilary and
Cassiodorus, graced indwelling in Augustine, and moral single-mindedness in
Theodoret.

At the antiphon's close, Hilary and Augustine hear bodily resurrection in the
strength God gives: Augustine explicitly invokes 1~Corinthians~15:43,
``sown in weakness, raised in power.'' Cassiodorus instead stresses patience
and firmness of faith on the road to eternal reward. The two readings should
not be flattened into a consensus. Both nevertheless forbid reading
``strength'' as domination by the gathered household. The divine power that
raises bodies or sustains fidelity is given to the people; it is not a license
to identify contemporary opponents as enemies whom worshippers may scatter.

\subsection*{Collect: mercy beyond merit and request\enspace\properrefs{Coll.}}

The prayer's grammar moves through three measures. God's abundance surpasses
the merits of suppliants, exceeds their requests, and gives beyond what prayer
itself presumes. Conscience supplies the negative petition---pardon what it
fears---while mercy supplies a positive good too great to originate in the
petition. Cummiskey's English shifts the final subject from ``what prayer does
not presume'' to blessings ``we dare not presume to ask.'' It remains the
identified historical translation, but the Latin is more exact about prayer's
own limit.

The entire collect-secret-postcommunion set is already joined in the Old
Gelasian, Book~III, VII, one of its then-unassigned Sunday Masses. The
Hadrianum Gregorian assigns it to the Twelfth Sunday after Pentecost; later
alignment places it at the Eleventh. This is transmission evidence, not an
originally fixed Sunday interpretation. The Collect was not located in the
Veronese collection, although the other two orations occur there separately.
Its oldest checked witness in this research is therefore the Gelasian, not a
conjectured Leonine source.

Aquinas supplies a doctrinal analogue without claiming to comment on the
prayer. In \emph{Summa theologiae} II--II, q.~83, a.~2, ad~3, he says that
God's liberality gives many things even when they have not been requested;
at q.~83, a.~16, he distinguishes a sinner's prayer heard through sheer mercy
from a demand of justice. Schuster's later liturgical commentary takes the
unpresumed gift toward perfected love and union with God. Together these
sources justify a restrained conclusion: mercy does not merely erase a
deficit but bestows a good beyond the petitioner's measure. They do not
identify exactly what temporal good will be given or guarantee felt
consolation.

\subsection*{Epistle: a transmitted Resurrection gospel\enspace\properrefs{Ep.}}

The lesson stands at the opening of Paul's long argument about resurrection.
Four verbs govern its first sentences: he preached, the Corinthians received,
they stand, and they are saved if they hold fast. Paul then handed on what he
received. Many scholars identify vv.~3--5 as a tradition earlier than the
letter; whatever the precise pre-Pauline form, the literary argument presents
the gospel as received testimony rather than Paul's invention. Death ``for
our sins,'' burial, Resurrection on the third day, fulfillment of Scripture,
and appearances belong together.

Irenaeus, \emph{Against Heresies} III.18.3, and Tertullian,
\emph{On the Resurrection of the Flesh} 48, use this sequence against
docetic or spiritualizing accounts: the Christ who truly died and was buried
rose in the flesh. Chrysostom, \emph{Homily 38 on First Corinthians}, treats
``delivered'' and ``received'' as the marks of transmission, burial as
confirmation of real death, and the appearances as cumulative witness. He
also sees Paul's self-description as deliberate humility: last appearance
does not mean a different gospel. Theodoret similarly emphasizes the
checkability of living witnesses and reads ``according to the Scriptures''
through Isaiah, the Psalms, and Jonah.

The appearance ``to the eleven'' differs from the Greek ``to the twelve.''
Ancient and medieval interpreters variously treat the larger number as the
corporate name of the apostolic body or regard the Vulgate's eleven as a
scribal or explanatory form. The Missal's \latin{undecim} controls this
publication; the textual history is reported without harmonization. The
appearances to James and ``all the apostles'' likewise extend beyond narrated
Gospel scenes. Aquinas notes more generally that the Gospels do not record
every listed appearance (\emph{Summa theologiae} III, q.~55, a.~3, ad~3).

Paul's \latin{tamquam abortivo} attracts genuinely different readings.
Chrysostom and Ambrosiaster stress conversion outside the ordinary apostolic
time; Theodoret treats the image as extreme self-abasement, like one not
entered among the living; Aquinas's \emph{Super I ad Corinthios} XV, lect.~1,
offers three aspects---untimely, violently brought forth, and not grown to due
measure. The work is a medieval \emph{reportatio vulgata}, not an autograph,
and the threefold analysis is reception rather than the one recoverable
meaning of Paul's metaphor. All preserve the contrast between undeserved
calling and real apostolic witness.

The last sentence demands the closest boundary control. The liturgical lesson
ends after ``his grace in me hath not been void.'' The canonical verse goes on
to Paul's abundant labor and ``the grace of God with me.'' Augustine,
\emph{On Grace and Free Will} 5.12, Gregory, \emph{Moralia in Iob} XVI.25,
and Aquinas use the complete sentence to articulate prevenient and cooperating
grace. Catholic doctrine receives that canonical development. Yet those
authors are not expounding the appointed ending alone, and their decisive
words were not proclaimed. The liturgical boundary concentrates attention on
grace transforming a persecutor and not becoming empty; it must not be used
to deny the labor that Scripture states one clause later.

\subsection*{Gradual: help, renewed flesh, and feared silence\enspace\properrefs{Grad.}}

The Gradual reverses the psalm's order: v.~7 is sung before v.~1. Its body
moves from a trusting heart to help, renewed flesh, and willing praise; the
versicle then returns to the cry that God not remain silent or depart. The
chant also differs from the Clementine at several points. It has
\latin{In Deo} where the biblical text has ``in him,'' a past
\latin{clamavi} where the Douay has ``will I cry,'' and
\latin{ne discedas a me}, a clause absent from that verse. No checked
patristic lemma carries all of this chant form.

Augustine, \emph{Enarratio in Psalmum} 27.7, reads
\latin{refloruit caro mea} with striking brevity: the flesh flourished again,
that is, rose again. Cassiodorus expands the botanical figure: the flesh that
first flowered in Christ's virginal conception flowers anew in the
Resurrection. Aquinas moves through Adam's lost innocence to Christ's risen
flesh and the members joined to him. This line of reception brings the
Gradual into direct proximity with the Epistle's Resurrection proclamation.

Theodoret gives the necessary dissent. In his exposition of Psalm~27 he reads
David's recovered vigor after the dangers associated with Saul, Doeg, and the
Ziphites; the flourishing is restored health and courage, not a direct
prophecy of Resurrection. The literal-historical and Christological readings
need not cancel one another, but the latter may not be presented as the only
ancient interpretation. The chant's plea against silence also remains a plea:
its later placement beside Christ's groan and word is liturgical relation, not
proof that the psalm originally described Mark's event.

\subsection*{Alleluia: joyful sound ordered to the true helper\enspace\properrefs{All.}}

Psalm~80 is a festival summons that quickly becomes a covenant summons to
listen: Israel is called from Egyptian bondage, forbidden strange gods, and
invited to receive God's provision. The Alleluia retains rejoicing, Jacob,
the psalm, and the harp, but drops the timbrel and psaltery and moves
``pleasant'' from the psaltery to the psalm. Because ancient commentary often
explains the omitted instruments, it illuminates the source context without
describing every word actually sung.

Augustine, \emph{Enarratio in Psalmum} 80.3--5, contrasts shouting in the
circus with rejoicing in God the helper. He reads the instruments as an
exchange between spiritual reception and embodied works of mercy: heavenly
preaching and earthly action must sound together. Cassiodorus develops that
moral register through almsgiving, fasting, and the union of divine words
with human acts. Theodoret remains closer to the historical festival and the
orchestra attested in Chronicles. That divergence is useful. Joy need not be
dematerialized into allegory, and musical celebration is not sufficient
without covenant hearing. In the Sunday sequence, the Alleluia's audible
praise precedes the command that ears and tongue be opened.

\subsection*{Gospel: \latin{Ephphetha} and the discipline of speech\enspace\properrefs{Gosp.}}

Mark situates Jesus on an indirect journey from Tyre by Sidon toward the Sea
of Galilee through the Decapolis. The route follows the Syrophoenician woman's
faith and keeps the narrative in or near Gentile territory; it precedes the
feeding of the four thousand. Matthew~15:29--31 gives a broad healing summary,
but no true parallel supplies this individual, the fingers, saliva, groan, or
Semitic command conventionally identified as Aramaic. A Hebrew classification
has also been argued. The 1962 incipit drops Mark's ``again'' and supplies Jesus'
name. Its controlling Vatican page prints both \latin{Exiens} and
\latin{finibus} without accents.

The sequence of acts resists reduction to one code. Others bring the man;
Jesus separates him from the crowd; fingers enter his ears; saliva touches
his tongue; Jesus looks heavenward and groans; one command opens; then the
man hears and speaks rightly. The bodily particularity dignifies the person
and the senses. A sacramental or moral reading begins from, and may never
cancel, this bodily good. Disability is neither personal guilt nor a convenient
name for unbelief.

Gregory the Great directly interprets vv.~32--35 in \emph{Homiliae in
Hiezechihelem} I.10.20. Fingers signify the Holy Spirit's gifts; saliva, the
wisdom of divine speech; the heavenward look, the source from which hearing
and speech are received; and the sigh, not necessity in Christ, but a summons
to direct our own groaning in prayer toward heaven. Gregory concludes that speaking rightly requires doing before
teaching. Under the 1960 Roman Office this paragraph is the third Matins
lesson for the Eleventh Sunday, after Mark~7:31--37. The liturgical fact makes
Gregory a privileged witness to received Sunday interpretation while not
turning his symbols into the literal sense of Mark.

Bede, \emph{In Marci evangelium expositio} II on 7:31--37, develops the
separation from the crowd as a first gate of healing, fingers as the Spirit,
saliva as the savor of wisdom, and the opened tongue as catechetical
confession. His account of Christ's sigh distinguishes the humanity that
looks upward from the divine authority of the effective word. The anonymous
early medieval Mark commentary transmitted under Jerome reads the healed
person corporately as humanity and connects the opened mouth with the Creed.
Its attribution must remain qualified; it is not Jerome's secure work.
Theophylact emphasizes both Christ's humility and the sanctifying efficacy of
his assumed flesh, while offering more than one explanation of the groan.

The Church's baptismal rites received the gestures in a different, ritual
register. Ambrose, \emph{De mysteriis} I.3--4, and the traditionally
Ambrosian \emph{De sacramentis} I.1.2--3, whose authorship has been disputed,
describe the opening rite of ears and nostrils through
\emph{Ephpheta}, so that the candidate may hear the mysteries and answer.
The 1952 Roman Ritual directs the priest's saliva-moistened thumb to the ears
and then the nostrils, never to the tongue; a 1944 decree permitted omission
of saliva for hygiene while retaining the touch and formulas. The rite is
preparatory and exorcistic, not the matter or form of Baptism.
Robert Bellarmine, \emph{De baptismo} I.24, later defends the received use of
saliva and the command. This is documented ritual reception, not evidence
that Mark narrated a later rite or that the Mass's Gospel itself performs it.

Saint Anthony of Padua's sermon on this Gospel, \emph{Dominica XII post
Pentecosten} §§3, 8--14, belongs to a different medieval Sunday numbering and a
different combination of propers. He reads the two ears as hearing the Old
and New Testaments, the fingers through the Spirit, and right speech as a life
whose confession is not contradicted by deeds. Saint Lawrence of Brindisi,
\emph{Dominicalia}, Sunday XI after Pentecost, develops
\latin{Bene omnia fecit} through creation and Isaiah~35, then treats the
gestures, baptism, hearing, and confession. His inherited equation of the
deaf man with unbelieving Israel is recorded only as historically harmful
reception; it is rejected as a synthesis. Catholic teaching in
\emph{Nostra aetate}~4 and the \emph{Catechism} §§839--840 forbids contempt
for the Jewish people and preserves God's irrevocable relation with Israel.

Augustine, \emph{De consensu evangelistarum} IV.4.5, concentrates on the
command not to tell. If forbidden witnesses publish eagerly, authorized
preachers should not be sluggish. Yet the command also guards Mark's story
from making publicity the measure of faith. Saint John Henry Newman later
notes that friends and the Church bring sufferers to Christ and that restored
speech serves confession. Saint Oscar Romero, preaching on 9 September 1979,
uses the concrete gestures to defend embodied human dignity and integral
liberation. Pope John Paul~II on 10 September 2000 reads \emph{Ephphatha} as
an opening to communion and intellectual mission. Pope Leo~XIV's general
audience of 30 July 2025 emphasizes Christ's healing intimacy, the overcoming
of isolating closure, and the need for the full paschal journey before speech
becomes mature confession. These later uses differ in genre and authority;
none may be stacked as if it were patristic consensus.

\subsection*{Offertory: the healed singer offers praise\enspace\properrefs{Off.}}

The Offertory takes two verses from an individual thanksgiving. Psalm~29
moves from rescue and healing through the memory of descent toward the pit,
the change of mourning into joy, and the vow of lasting praise. The chant
omits \latin{Deus meus} from v.~3. Its position after Mark creates an audible
liturgical correspondence---``I cried to thee, and thou hast healed me''---but
does not establish that the psalm was selected as an original commentary on
this one miracle.

Augustine's two expositions read the psalm through Christ and the Church.
The dedication of the house becomes Resurrection: Christ's healed flesh is
not recovery from sin or disease but the bearing and overcoming of human
mortality. Cassiodorus closely develops the same paradox, identifying the
house with Christ's body and dedication with Resurrection. Aquinas joins that
Christological house to the Church still being built in conversion and
awaiting dedication. Theodoret independently interprets the title as the
renewal of human nature in Christ, while allowing the verse's immediate sense
to fit Hezekiah's healing.

This tradition draws the Offertory toward the Epistle's risen body and the
Postcommunion's salvation of body and mind. Its guardrail is Christology:
``healed'' cannot imply moral or physical defect in Christ. Nor does chanting
the verse promise an identical temporal cure to every Christian. The
resurrection reading places bodily healing within the promised destruction of
death.

\subsection*{Secret: accepted service becomes help in fragility\enspace\properrefs{Sec.}}

The prayer asks God to look favorably upon the Church's service so that what
she offers may be an accepted gift and assistance for her fragility. The
offering is real, but neither its acceptance nor its fruit is self-produced.
The repeated \latin{sit} keeps both results in the form of petition. The
Trinity Preface follows by rubric; it is not another variable proper and is
therefore not reproduced here.

The prayer is older than its Sunday setting. The Veronese sacramentary places
an early form within a Mass in time of drought; it has
\latin{ad nostram servitutem} and the dative
\latin{fragilitati}. The 1570 Missal still carries \latin{ad}, while the 1604
Vatican typical edition has already dropped it and approaches the 1962
syntax. The Old Gelasian unites the Secret with this Sunday's Collect and
Postcommunion but lacks the second \latin{et}. These facts establish a long
transmission and real variation, not the prayer's author or original
theological program.

Schuster's identification of \latin{servitus} with the Church's
\emph{leitourgia}, her priestly ministry, offers an illuminating reception.
The prayer itself is broader than an instruction about the ordained minister:
the whole liturgical ``we'' asks that offered service be acceptable and
medicinal. Human weakness is neither denied nor treated as a reason not to
offer. It is brought into a relation where service received by God becomes
help for the servants.

\subsection*{Communion: substance, firstfruits, and a guarded promise\enspace\properrefs{Comm.}}

The antiphon compresses Proverbs~3:9--10. It retains honoring the Lord from
one's substance and firstfruits, then the full barns and overflowing presses;
it omits the Clementine's ``all,'' the imperative ``give him,'' and one
possessive. In context this saying belongs to a father's wisdom instruction:
trust in the Lord rather than one's prudence, accept discipline, prize wisdom,
show mercy, and do not postpone a neighbor's due. The imagery therefore joins
worship and economic goods inside a moral formation; it is not a detachable
formula for obtaining wealth.

Bede, \emph{In Proverbia Salomonis} I on 3:9--10, gives three levels. His
literal reading requires generosity from goods that justly belong to the giver
and directs the act to God, not to the giver's reputation. Spiritually, the firstfruits include
every good work and grace credited to God rather than one's independent
power. Eschatologically, filled barns and presses figure heavenly satiety.
Rabanus transmits this account. Salvian of Marseilles,
\emph{Ad Ecclesiam} I, remains more sharply economic: because all possessions
are received from God, returning them to him is justice as well as devotion.

The interpretations converge on received possession, gratitude, and the
neighbor, but they do not make poverty a sign of failed faith. The Lord whom
Communion honors is the crucified and risen Christ received sacramentally;
Christian tradition must read material abundance beside almsgiving, the
danger of avarice, and the eschatological good. Schuster connects the
harvest-time antiphon with the old blessing of firstfruits, a plausible
liturgical reception rather than proof of the Communion's first assignment.

\subsection*{Postcommunion: one salvation of mind and body\enspace\properrefs{Postcomm.}}

The assembly asks to experience, through reception of the sacrament, help of
mind and body, so that saved in both it may glory in the fullness of heavenly
remedy. Cummiskey reverses the Latin order to ``body and mind'' and renders
\latin{plenitudine gloriemur} as enjoying the remedy's full effect. Those are
features of the quoted nineteenth-century witness, not replacements for the
Latin's movement from perception through aid and salvation to glory.

An early form occurs in the Veronese collection and includes a
\latin{de} before \latin{caelestis remedii plenitudine}; the Gelasian and
Gregorian transmit the prayer within the same set as the Collect and Secret.
The prayer's anthropology is positive and eschatological. Salvation concerns
the embodied person; sacramental remedy is heavenly; and its fullness exceeds
the present moment. It neither promises that Communion removes every illness
now nor relegates the body to a temporary shell.

Schuster makes the resurrection horizon explicit: bodily contact with
Christ's sacramental Body strengthens the Christian toward the body's final
raising. The connection is later liturgical theology, but it fits the day's
source field with unusual precision. Paul's risen Christ, the Gradual's
flourishing flesh, Mark's restored senses, the Offertory's healing, and the
Postcommunion's salvation of mind and body all resist a Christianity reduced
to disembodied information. The last petition also prevents triumphalism:
those who have received still confess their need for help.
